% GENERATED by md2tex.py from sections/03_rank_law.md - edit the .md, not this file
\section{The Rank Law}

In the encoder family, the task's representation theory sets the
state's dimensionality.

\subsection{Recruited Rank Tracks Minimal Faithful Dimension}

Figure 1 shows recruited restricted effective rank against $d_{\min}$
for the five groups, per seed. The per-group means are 1.877 ($S_3$),
2.852 ($S_4$), 2.832 ($A_5$), 3.591 ($S_5$), and 4.736 ($A_6$) against
$d_{\min} = 2, 3, 3, 4, 5$ \% evidence: R1. Spearman correlation
across the family is $\rho = 0.9747$ \% evidence: R1. This number
is the maximum the family permits: the $S_4$/$A_5$ tie at $d_{\min}=3$
caps the achievable midrank-adjusted $\rho$ at exactly 0.9747, and the
observed ordering attains that cap \% evidence: R1. Because the
family has only five members, we report the exact permutation null
rather than an asymptotic one: under the null, $P(\rho \ge 0.8) = 8/120 \approx 6.67\%$
\% evidence: R1. The correlation is therefore suggestive on its
own; the band criterion and the causal test below do the deciding. The
tracking also holds cell by cell: all 19 unconstrained-arm cells in the
sweep land inside the pre-registered $[0.7, 1.3] \cdot d_{\min}$ band
\% evidence: R1.

\subsection{The Designed Dissociation: $S_4$ versus $A_5$}

$S_4$ is solvable; $A_5$ is the smallest non-solvable group. Both have
$d_{\min}=3$. If recruited rank were driven by computational difficulty
in the Krohn-Rhodes sense, the pair should separate; if by
representation dimension, the pair should land together. We
pre-registered a two-one-sided-tests equivalence procedure (Welch, $n=5$
per group, margin $\pm 0.5$ rank-units, half the ladder's adjacent-rung
spacing). The result is a declaration of equivalence: mean difference
0.0194 rank-units (2.852 versus 2.832), standard error 0.0368, $df =
7.83$, with both one-sided statistics ($t_1 = 13.06$, $t_2 = 14.12$)
exceeding the critical value 1.865 by a factor of seven
\% evidence: R2. Recruited rank tracks the representation
dimension of the task; the group's algebraic complexity class
contributes nothing detectable.

\subsection{An Instrument Defect, Its Diagnosis, and the Repaired Causal Test}

An early force-rank sweep produced a mechanically valid FALSIFY reading
that a target-embedding defect later voided, and the defect changes
what the causal test means. The training target had been embedded as
$\rho(g) \oplus I_{d_{state}-d_{\min}}$: an eye-padded block matrix of
full rank $d_{state}$ with all singular values equal to one. For such a
target, a rank-$k$ model's optimal direct cosine is $\sqrt{k/d_{state}}$
independent of which subspace it spends its rank on, and the harvest
matched this prediction: 36 of 39 force-rank cells sat within 0.07 of
the predicted ceiling, one further cell sat marginally outside at
0.072, and the two remaining outliers (0.15 and 0.17, both in one
group's below-$d_{\min}$ arm) reflect an additional optimization
shortfall \% evidence: R1b. Rank-capped models were
buying cosine from the identity padding, so the causal boundary at
$d_{\min}$ was never tested. The fix replaced the padding with zeros,
making the target rank exactly $d_{\min}$, and the decisional metric was
pre-registered before the re-run as the Procrustes-aligned crosscheck
($\mathrm{xrec90}$), because the scale-only primary metric reads cosine
0.01 to 0.22 on a flawless oracle when the target's rho-block
eigenvalues are degenerate \% evidence: R1b.

\subsection{Recovery Is a Step Function at Exactly $d_{\min}$}

Figure 2 shows the repaired force-rank staircase. At $k = d_{\min}-1$,
held-out recovery $\mathrm{xrec90}$ is 0.000 in all five groups
\% evidence: R3: not degraded, zero. At $k = d_{\min}$, recovery
returns and clears the pre-registered bar of 0.9 times each group's own
unconstrained anchor in four of five groups directly ($S_4$ 0.800
against bar 0.585; $A_5$ 0.700 against 0.630; $S_5$ 0.600 against 0.450;
$A_6$ 0.650 against 0.585) \% evidence: R3. $S_3$ landed within
the pre-stated marginality window (0.450 against bar 0.495) and was
routed to its pre-registered seed extension: across four independent
seeds the seed-mean at $k = d_{\min}$ is 0.5625, clearing the bar, and
$k = d_{\min}-1$ reads exactly 0.000 in all four seeds
\% evidence: R3. The pre-registered confirmation criterion (a
razor step at $d_{\min}$ in at least four of five groups, including at
least one of the marquee pair) is met with both marquee groups
confirming \% evidence: R3.

Together the three readings close the loop in both directions:
$d_{\min}$ is recruited when available and is causally necessary and
sufficient when forced. The substrate is small and every reading
reproduces from the archived artifacts (Appendix B).

\begin{figure}[t]
\centering
\includegraphics[width=0.62\linewidth]{fig1_rank_vs_dmin.pdf}
\caption{Recruited restricted effective rank versus minimal faithful representation dimension $d_{\min}$ for five permutation groups ($S_3$, $S_4$, $A_5$, $S_5$, $A_6$; unconstrained arms, per-seed points, group means marked). The shaded region is the pre-registered $[0.7, 1.3] \cdot d_{\min}$ band; all 19 cells land inside it. Spearman $\rho = 0.9747$ equals the tie-capped maximum for this family. The $S_4$/$A_5$ pair (both $d_{\min}=3$, opposite solvability) lands together; the equivalence test in Section 3.2 declares them equivalent within $\pm 0.5$ rank-units.}
\label{fig:rank-vs-dmin}
\end{figure}

\begin{figure}[t]
\centering
\includegraphics[width=\linewidth]{fig2_forcerank_staircase.pdf}
\caption{Held-out recovery ($\mathrm{xrec90}$, Procrustes-aligned cosine threshold 0.9) versus forced train-time rank $k$, per group, on the zero-padded (repaired) target. Recovery is 0.000 at $k = d_{\min}-1$ in every group and returns at $k = d_{\min}$; dashed lines mark each group's bar of 0.9 times its own unconstrained anchor. The $S_3$ panel overlays all four extension seeds: every seed reads 0.000 at $k = d_{\min}-1$, and the seed-mean at $k = d_{\min}$ (0.5625) clears the bar (0.495).}
\label{fig:forcerank}
\end{figure}


